\subsection{Five deletions for two missing-edge configurations}
The cover-gain-three cases depend on the shape of the missing edges. Two shapes permit an additional exclusion.

\begin{proposition}[A path or triangle and a disjoint edge]\label{prop:pathfive}
Let $G$ have no four-cycle, and let a positive variable-input binary circuit compute $P\in\Sup(G)$ with nonempty quadratic support $H$. Suppose the graph formed by its four missing edges is $P_4\sqcup K_2$ or $K_3\sqcup K_2$, where $P_4$ is a path with three edges. If $d=3$ and the extracted quadratic circuit is optimal, then five deletions are impossible.
\end{proposition}
\begin{proof}
Let $J$ be the missing-edge graph, $t=\tau(H)$, and $Z$ the assigned minimum cover of the optimal quadratic frame. Both possible $J$ have cover number three, and deleting any one of their vertices lowers that number to two. If such a vertex belonged to $Z$, adjoining a cover of the remaining missing graph would cover $G$ with $t+2$ vertices. Hence $Z$ avoids every missing endpoint. On $T=V(G)\setminus Z$, the graph is precisely $J$ with possible isolated vertices.

Apply full-carrier erasure before restriction. No private gate disappears, by the earlier five-deletion lower bound for four missing witnesses. Every nonzero nonconstant restricted direct context of a visible cut is not contained in a minimum cover of $J$. Indeed, it perturbs a literal assigned center $h\in Z$ and contains $N_H(h)\cap T$. If it lay in a minimum $J$-cover $U$, replacing $h$ in $Z$ by $U$ would cover $H$ and all missing edges with $t+2$ vertices. No missing edge is incident to $h$, since $Z$ avoids their endpoints. Zero and scalar contexts remain allowed. A union of permitted nonzero linear contexts is still not contained in any minimum $J$-cover.

A second restriction comes directly from four-cycle-freeness. If $u,w,v$ form a two-edge path of missing edges, no old linear boundary clause may contain both $u$ and $v$. Such a surviving clause has a quadratic product above it and therefore a common old neighbor $h$ of $u,v$. Squarefreeness excludes $h=u,v$; $h=w$ would make a missing edge old; and any other $h$ produces the four-cycle $u,w,v,h$. The same restriction holds for one individual nonconstant visible-cut context $D$, because its nonzero old low part $A$ satisfies $AD\subseteq P_2$ in support. It need not hold for a union of contexts from different cuts.

Identify all isolated vertices with one variable $z$. This preserves the target, the minimum-cover context condition, and the missing-path restrictions on the genuine endpoints. We now consider the number of deleted products.

With at most two products, full-carrier expansion gives the general two-product form
\[
 g_i=A_i(x,g_0,\ldots,g_{i-1})B_i(x,g_0,\ldots,g_{i-1}),\qquad
 F=\bigvee_i g_iL_i\vee\!\bigvee_{i<j:\delta_{ij}=1}g_ig_j.
\]
The operand clauses are unrestricted positive disjunctions, and the affine contexts obey only the minimum-cover condition. A direct context is a union of visible-cut contexts; diagonal terms become constant-one contexts by Boolean idempotence. In this class we impose no missing-path restriction on the operands or the merged contexts. Private additions can create new linear sums, and different cuts can have different old neighbors. The finite certificate below excludes this larger class for both targets.

Suppose there are three products and two additions. If both additions are visible cuts selecting the same high product, all three products are ancestral to it, and it is a conjunction of at most four old clauses. One merged affine context gives at most five CNF clauses, whereas either target requires six. If the two highs are distinct, no private addition remains. Each operand is an old linear clause, possibly with inserted highs, or an earlier product alone. We enlarge this class by allowing every earlier product in any operand disjunction. Its old variable part satisfies the missing-path restriction, and each of the two individual direct contexts satisfies both restrictions. Only two product indices can appear directly or in output pair terms. The second finite class below excludes this possibility.

If only one addition is visible, the high-input computation is a four-node Boolean DAG with three AND gates and one OR gate, over old linear boundary clauses. Each old clause obeys the missing-path restriction. The private OR itself may create a forbidden pair, so it is represented as an explicit gate with unrestricted output. Each input is an old clause or an earlier node, and the final affine context obeys both restrictions. The third finite class below excludes this case, including every operation order and all sharing.

Finally suppose four products are deleted. The sole deleted addition would contradict Theorem~\ref{thm:purecut}, since $d=3$ requires $D\ge2^3-1=7$. At least one deleted addition is necessary, so these product counts exhaust the possibilities.

Here are the complete finite verifications used in the first three cases. For the first two classes, choose each variable and each permitted earlier product in an operand by a Boolean selector, and select affine contexts and pair terms similarly. The minimum-cover condition says that a nonzero nonconstant context has a selected variable outside each minimum cover. In the second class, prohibit each missing-path pair in old variable parts and individual contexts; three activity bits dominate every context and pair selector, with at most two bits true. In the third class, one-hot operand modes choose an old clause or an earlier node, and operation bits select exactly one OR among four gates. Its old clauses and final context obey the stated restrictions.

Exact Tseitin equivalences impose truth on the four edge assignments with $z=0$ and falsity on the six maximal independent assignments with $z=1$. Monotonicity makes these ten tests complete. Empty clauses, zero or constant-one contexts, repeated operands, unused nodes and restricted coincidences are permitted. The resulting counts are
\[
\begin{array}{c|l|r|r|r}
J&\text{class}&\text{variables}&\text{clauses}&\text{proof lines}\\\hline
P_4\sqcup K_2&\text{two products}&632&2157&493\\
P_4\sqcup K_2&\text{two visible cuts}&992&3442&2210\\
P_4\sqcup K_2&\text{one private OR}&1130&3838&1678\\
K_3\sqcup K_2&\text{two products}&566&1913&524\\
K_3\sqcup K_2&\text{two visible cuts}&893&3085&2873\\
K_3\sqcup K_2&\text{one private OR}&1031&3486&1093.
\end{array}
\]
All six CNFs were independently regenerated byte for byte. Complete DRUP refutations were independently verified in RUP-only mode with zero RAT steps; the checks took between $0.06$ and $0.27$ seconds each. The program \emph{Five-deletion obstruction}, exact source bindings, compressed proofs and checker logs accompany the workbench. The certified classes contain all cases above, completing the proof.
\end{proof}

At cover gain three with five deletions and an optimal frame, the missing graph therefore cannot be $P_4\sqcup K_2$ or $K_3\sqcup K_2$. The shape $P_3\sqcup K_2\sqcup K_2$, and configurations in which a minimum old cover meets a missing endpoint, require further arguments. A center exchange in the latter case can uncover a missing edge incident to the removed center.
